\section{Lean formalization architecture}

\subsection{Public theorem shape}

The implementation should expose one structure for a compiled fixed formula.
A simplified signature is:

\begin{lstlisting}[language={},caption={Core package shape}]
structure CompiledFormula (n : Nat) where
  meaning : Fm PTm n
  value positiveBound negativeBound : PTm n
  positiveCheck negativeCheck : PTm (n + 1)
  canonicalPositive canonicalNegative : PTm n

  valueBoolean : [] |- [PRAAx] value = 0 or value = 1
  valueCorrect : [] |- [PRAAx] (value = 1 <-> meaning)
  positiveSound : [positiveCheck = 1] |- [PRAAx] meaning.shift
  negativeSound : [negativeCheck = 1] |- [PRAAx] not meaning.shift
  positiveCodeBound : [positiveCheck = 1] |- [PRAAx]
    var 0 < positiveBound.shift
  negativeCodeBound : [negativeCheck = 1] |- [PRAAx]
    var 0 < negativeBound.shift
  canonicalPositiveAccepted : [] |- [PRAAx]
    value = 1 -> positiveCheck[canonicalPositive] = 1
  canonicalNegativeAccepted : [] |- [PRAAx]
    value = 0 -> negativeCheck[canonicalNegative] = 1
\end{lstlisting}

The actual source in \code{lean/STRIVE/CompiledFormula.lean} uses the repository
syntax and scope operations.  The key point is that proof fields are ordinary
object derivations.  A semantic theorem such as
\code{evaluate formula env = true} is not accepted as a substitute.

\subsection{Module dependency order}

The recommended dependency chain is:

\begin{enumerate}
\item \code{STRIVE.Kernel}: boundary, specialization, simultaneous cut,
certificate sealing.  This file has no placeholders.
\item \code{Core.MixedRadix}: prefix products, packed sums, digit bounds,
decoding, and extensionality.
\item \code{Core.BooleanFolds}: exact all/any equations, prefix sums/products,
and disjoint-owner search.
\item \code{STRIVE.CompiledFormula}: package structure, atom, conjunction, and
disjunction constructors.
\item Bounded existential and universal constructors.
\item Generated formula modules and regression theorems.
\item Contract graph, demand closure, trace/viability, and observer modules.
\item Foundation/open-sequent and native exports.
\end{enumerate}

This order keeps the trusted surface small.  The mixed-radix module is useful
independently and can be reviewed before any generated compiler code.

\subsection{Mixed-radix proof path}

The first substantial Lean milestone should be
\code{IPRADerivations/Core/MixedRadix.lean}.  Reuse the exact current theorem
names where possible:

\begin{itemize}
\item \code{prod2\_zero}, \code{prod2\_succ}, \code{sum2\_zero}, and
\code{sum2\_succ} for prefix/table recurrences;
\item \code{div\_mod} and \code{div\_mod\_unique} for digit extraction;
\item \code{sum2\_congr\_bounded} for extensionality;
\item order and semiring lemmas for the append interval;
\item \code{ind} or \code{PRA.qfInd} for code bounds and decoding.
\end{itemize}

A useful local normal form is
\[
 \pack(N+1)=\pack(N)+a_N\pref(N),
 \qquad
 \pref(N+1)=\pref(N)B_N.
\]
Prove these equations at arbitrary argument terms first.  Then the induction
steps become arithmetic rewrites plus one interval lemma.

\subsection{Bounded folds}

The compiler needs two families of finite-fold lemmas.

For \(\mathsf{all}\):
\begin{align*}
 \mathsf{all}(0)&=1,\\
 \mathsf{all}(N+1)&=\mathsf{all}(N)\cdot p(N),\\
 \mathsf{all}(N)=1&\to p(i)=1\quad(i<N),\\
 (\forall i<N\;p(i)=1)&\to\mathsf{all}(N)=1.
\end{align*}
For \(\mathsf{any}\), use the dual recurrence and least-success theorem.  The
current search library finds the least zero of a body; a least one is obtained
by searching the Boolean complement or an equality-to-one test.

These lemmas should be object derivations parameterized by concrete body terms,
not host propositions over arbitrary functions.

\subsection{Constructor implementation}

Each syntax constructor is a Lean definition returning a
\code{CompiledFormula} record.

\paragraph{Atom.}
Fill every field by the reflection package and numeral simplification.

\paragraph{Conjunction and disjunction.}
Define code bounds by term addition/multiplication.  Use the two-digit
mixed-radix special case or pairing lemmas for product codes, and a two-interval
disjoint sum for choices.  Prove checker soundness directly with natural-
deduction introduction/elimination.

\paragraph{Quantifiers.}
Generate a concrete child-bound body in one extra scope.  Prefix sums and
products are \(PR\) symbols.  The checker's finite loop is a specialized
primitive recursion.  Soundness is a quantifier-free induction on the bound;
child soundness is instantiated at the current numerical index by
\code{Prf.subst}.

A practical first implementation may use a slightly redundant code format if
that shortens the proof.  The exact tight heterogeneous bounds can then replace
it behind the same package interface.

\subsection{Metalevel generator}

The generator reads a restricted description of a fixed formula and emits a
Lean module.  It performs:

\begin{enumerate}
\item scope checking and alpha-renaming;
\item allocation of concrete \(PR\) symbols for every value, bound, checker,
prefix sum, and prefix product;
\item constructor assembly;
\item theorem-name generation and a machine-readable ledger;
\item optional reference fixtures for the Python model.
\end{enumerate}

The generator is not in the trusted base.  Its output must elaborate without
\code{sorry}.  A small formula should generate readable code before any effort
is spent on optimization.

\subsection{Open-sequent integration}

The contract compiler should retain a list
\code{residuals : List (Fm PTm n)} and a proof
\[
 \code{residuals |- [PRAAx] goal}.
\]
Closing a graph is exactly \code{proof.trans supply}.  Specializing a query is
exactly \code{proof.subst sigma}.  A proof-local certificate is exactly
\code{Prf.exE producer consumer}.  These structural operations are already
small enough to audit manually.

Only after this ordinary endpoint is complete should a client invoke the
open-sequent bridge.  The transported theorem must preserve the free-variable
assignment and every residual hypothesis.  A native scalar export is similarly
a final corollary, not the main implementation.

\subsection{Tests and audits}

A completed module should pass four independent checks.

\begin{enumerate}
\item \textbf{Equation tests}: numeral examples for every generated \(PR\)
symbol, especially zero bounds and empty products.
\item \textbf{Cross-model tests}: compare generated evaluator, bounds, codes,
and checker with \path{strive_reference.py} on small inputs.
\item \textbf{Proof tests}: instantiate public \(\Prf\;\PRAAx\) theorems at
small numerals and build the corresponding check modules.
\item \textbf{Axiom audit}: ensure no declaration depends on \code{sorryAx},
classical axioms, or an accidental semantic theorem.
\end{enumerate}

The Python exhaustive tests in this packet check every code in both signed code
spaces for the two principal examples.  They establish implementation
consistency, not formal soundness; the Lean derivations remain the authoritative
goal.

\subsection{Estimated formalization milestones}

The dependency structure suggests independently useful merge points:

\begin{enumerate}
\item mixed-radix arithmetic and disjoint sums;
\item quantifier-free Boolean folds;
\item binary signed formulas;
\item one bounded existential;
\item one bounded universal;
\item generated alternating examples;
\item contract escape modes;
\item demand closure and trace integration.
\end{enumerate}

Each milestone has mathematical value even if later generator work is delayed.
The detailed command-level checklist is in \code{lean/FORMALIZATION\_PLAN.md}.
